\documentclass[11pt,a4paper]{article}
\usepackage{fontspec,amsmath,amssymb,amsthm,microtype}
\setmainfont{lmroman10-regular.otf}[BoldFont=lmroman10-bold.otf,ItalicFont=lmroman10-italic.otf,BoldItalicFont=lmroman10-bolditalic.otf]
\usepackage{xcolor}
\usepackage[most]{tcolorbox}
\usepackage{mathrsfs}
\providecommand{\lt}{<}
\providecommand{\gt}{>}
\setmonofont{DejaVuSansMono.ttf}[BoldFont=DejaVuSansMono-Bold.ttf,ItalicFont=DejaVuSansMono-Oblique.ttf,BoldItalicFont=DejaVuSansMono-BoldOblique.ttf,Scale=0.86]
\newfontfamily\leanfallback{STIXTwoMath-Regular.otf}[Scale=0.86]
\definecolor{leanrule}{gray}{0.62}
\definecolor{leanback}{gray}{0.962}
\definecolor{leantitle}{gray}{0.22}
\definecolor{leanmuted}{gray}{0.38}
\newtcolorbox{leanbox}[1]{enhanced,breakable,colback=leanback,colframe=leanrule,
  boxrule=0.4pt,arc=0.6mm,left=2.4mm,right=2.4mm,top=1.8mm,bottom=1.8mm,
  before skip=0.8em,after skip=1em,fonttitle=\small\bfseries,coltitle=leantitle,
  colbacktitle=leanback,titlerule=0.3pt,title={#1}}
\newlength{\leancharwidth}
\newenvironment{leancode}{\par\smallskip\begingroup\ttfamily\footnotesize\raggedright
  \settowidth{\leancharwidth}{M}\setlength{\parskip}{0pt}\setlength{\parindent}{0pt}
  \frenchspacing}{\par\endgroup\smallskip}
\newcommand{\leanline}[2]{\hangindent=\dimexpr#1\leancharwidth+4\leancharwidth\relax
  \hangafter=1\noindent\hspace*{#1\leancharwidth}#2\par}
\newcommand{\leansource}[1]{\par{\footnotesize\color{leanmuted}#1}\par}
\newcommand{\leanlabel}[1]{\par\smallskip\noindent{\small\bfseries #1}\par}
\usepackage[margin=25mm]{geometry}
\usepackage[unicode,hidelinks]{hyperref}
\hypersetup{pdfauthor={LangevinLean Project}}
\setlength{\emergencystretch}{3em}
\setlength{\parskip}{0.4em}
\setlength{\parindent}{0pt}
\allowdisplaybreaks[2]
\raggedbottom
\title{\texorpdfstring{The hypotheses on which the main results rest}{The hypotheses on which the main results rest}}
\author{LangevinLean Project}
\date{}
\begin{document}\maketitle
Draft mathematical exposition; these notes may contain mistakes.


The principal statements of this report are proved from explicitly stated hypotheses only. For the potential: \(U\) is differentiable on \(\mathbb R^{d_\theta}\times\mathbb R^{d_x}\), \(m\)-strongly convex with \(m>0\), and \(\nabla U\) is \(L\)-Lipschitz, which forces \(m\le L\) whenever \(d_\theta+d_x>0\) and is assumed explicitly where it is used. For the particle system: \(N\ge1\). For the dynamics and the algorithms: a positive friction \(\gamma\) with \(\gamma^2\ge m+L\) for the diffusion; in addition a positive step size \(h\) with \(4\gamma Lh\le m\) for KIPLMC1; and \(\gamma^2\ge4L\) together with a positive step size satisfying \(33\gamma^3h\le m\) for KIPLMC2. For the initialization: an initial position law with finite second moment, and an initial velocity distributed as \(\mathcal N(0,N^{-1}I)\) in scaled coordinates, independent of the initial position.




All further ingredients are constructed or proved rather than assumed: the Brownian motion driving the diffusion, the solution of the particle equation, the stationarity of the Gibbs phase law, the correlated Gaussian innovations of KIPLMC1 and the independent thermostat noises of KIPLMC2, the Markov laws of both chains, the contraction of the OBABO step, and the finiteness of every Wasserstein distance that is converted to a real number. In particular no discrete invariant law of the OBABO chain and no external bias estimate is presupposed. The Brownian motion is constructed as a centered Gaussian process with the Brownian covariance whose values are continuous in mean square; the arguments use only its finite-dimensional laws and this mean-square continuity, so no modification with continuous sample paths is required.


\end{document}
